\begin{lemma}[Multifold sums of zero powers]
  \label{lem:multifold-sum-of-zero-powers}
  For non-negative integers $r,n$
  \begin{align*}
    \KnuthRFoldSum{r}{n}{0} = \binom{r+n-1}{r}.
  \end{align*}
  \begin{proof}
    \leavevmode\par
    \begin{enumerate}
      \item Let $r=0$, then we have $\KnuthRFoldSum{0}{n}{0} = 1$,
        by definition~\eqref{def:multifold-sums-of-powers-recurrence}.
      \item Let $r=1$, then we have
        $\KnuthRFoldSum{1}{n}{0}
        = \sum_{k=1}^{n} 1 = \binom{n}{1}$.
      \item Let $r=2$, then we have
        $\KnuthRFoldSum{2}{n}{0}
        = \sum_{k=1}^{n} \binom{k}{1}
        = \binom{n+1}{2}$.
      \item Let $r=3$, then we have
        $\KnuthRFoldSum{3}{n}{0}
        = \sum_{k=1}^{n} \binom{k+1}{2}
        = \binom{n+2}{3}$.
      \item Similarly, for arbitrary integer $r$, by hockey stick identity
        $\sum_{k=1}^{n} \binom{k}{r} = \binom{n+1}{r+1}$,
        yields
        $\KnuthRFoldSum{r}{n}{0}
        =  \sum_{k=1}^{n} \KnuthRFoldSum{r-1}{k}{0}
        =  \sum_{k=1}^{n} \binom{k+r-2}{r-1}
        = \binom{r+n-1}{r}$.
    \end{enumerate}
    This completes the proof.
  \end{proof}
\end{lemma}
Thus, by Lemma~\eqref{lem:multifold-sum-of-zero-powers}, and
by Theorem~\eqref{thm:multifold-sums-of-powers} yields
\begin{proposition}[Multifold binomial sums of powers]
  \label{prop:multifold-binomial-sums-of-powers}
  For non-negative integers $n,m$, and an arbitrary integer $t$,
  \begin{align*}
    \KnuthRFoldSum{r}{n}{m}
    &= \frac{1}{2} \sum_{k=0}^{m}
    \Bigg \{
      \binom{n-t+\frac{k}{2}+r}{k+r} + \binom{n-t+\frac{k}{2}+r-1}{k+r} \\
      &-
      \sum_{s=0}^{r-1}
      \left[
        \binom{r-s-t+\frac{k}{2}}{k+r-s} + \binom{r-1-s-t+\frac{k}{2}}{k+r-s}
      \right] \binom{s+n-1}{s}
    \Bigg \} \delta^{k} t^m.
  \end{align*}
\end{proposition}
For example,
\begin{example}[Examples of multifold binomial sums]
  For integer $0 \leq t \leq 2$
  \begin{align*}
    t=0:\quad
    \KnuthRFoldSum{r}{n}{m}
    &=
    \tfrac12
    \tsum_{k=0}^{m}
    \Bigg\{
      \tbinom{n+\tfrac{k}{2}+r}{k+r}
      +
      \tbinom{n+\tfrac{k}{2}+r-1}{k+r}
      \\
      &\qquad-
      \tsum_{s=0}^{r-1}
      \Big[
        \tbinom{r-s+\tfrac{k}{2}}{k+r-s}
        +
        \tbinom{r-1-s+\tfrac{k}{2}}{k+r-s}
      \Big]
      \tbinom{s+n-1}{s}
    \Bigg\}
    \,\delta^k 0^m ,
    \\
    %
    t=1:\quad
    \KnuthRFoldSum{r}{n}{m}
    &=
    \tfrac12
    \tsum_{k=0}^{m}
    \Bigg\{
      \tbinom{n-1+\tfrac{k}{2}+r}{k+r}
      +
      \tbinom{n-1+\tfrac{k}{2}+r-1}{k+r}
      \\
      &\qquad-
      \tsum_{s=0}^{r-1}
      \Big[
        \tbinom{r-s-1+\tfrac{k}{2}}{k+r-s}
        +
        \tbinom{r-2-s+\tfrac{k}{2}}{k+r-s}
      \Big]
      \tbinom{s+n-1}{s}
    \Bigg\}
    \,\delta^k 1^m ,
    \\
    %
    t=2:\quad
    \KnuthRFoldSum{r}{n}{m}
    &=
    \tfrac12
    \tsum_{k=0}^{m}
    \Bigg\{
      \tbinom{n-2+\tfrac{k}{2}+r}{k+r}
      +
      \tbinom{n-2+\tfrac{k}{2}+r-1}{k+r}
      \\
      &\qquad-
      \tsum_{s=0}^{r-1}
      \Big[
        \tbinom{r-s-2+\tfrac{k}{2}}{k+r-s}
        +
        \tbinom{r-3-s+\tfrac{k}{2}}{k+r-s}
      \Big]
      \tbinom{s+n-1}{s}
    \Bigg\}
    \,\delta^k 2^m .
  \end{align*}
\end{example}

By setting $t \rightarrow -t$ we get negated version
of Proposition~\eqref{prop:multifold-binomial-sums-of-powers}
\begin{proposition}[Multifold binomial sums of powers negated]
  \label{prop:multifold-binomial-sums-of-powers-negated}
  For non-negative integers $n,m$, and an arbitrary integer $t$,
  \begin{align*}
    \KnuthRFoldSum{r}{n}{m}
    &= \frac{1}{2} \sum_{k=0}^{m} (-1)^{m+k}
    \Bigg \{
      \binom{n+t+\frac{k}{2}+r}{k+r} + \binom{n+t+\frac{k}{2}+r-1}{k+r} \\
      &-
      \sum_{s=0}^{r-1}
      \left[
        \binom{r-s+t+\frac{k}{2}}{k+r-s} + \binom{r-1-s+t+\frac{k}{2}}{k+r-s}
      \right] \binom{s+n-1}{s}
    \Bigg \} \delta^{k} t^m.
  \end{align*}
  \begin{proof}
    By setting $t \rightarrow -t$ into
    Proposition~\eqref{prop:multifold-binomial-sums-of-powers}, and by
    $\delta^k (-t)^m = (-1)^{m+k} \delta^{k} t^{m}$.
  \end{proof}
\end{proposition}
There are indeed many (infinitely) variations of such formulas as above,
at least for every integer $t$.
However, this entropy can be even wilder, by recalling the fundamental properties
of binomial coefficients, such as
\begin{align*}
  \binom{-x}{k} &= (-1)^{k} \binom{x+k-1}{k}, \\
  \binom{x}{k}  &= \binom{x}{x-k}.
\end{align*}
